\section{Open questions and next steps}

The following five questions are genuinely open after this replication.
None is answered by the paper as published, and none is answered by
the on-disk artifacts produced here. Each is paired with a concrete,
free-endpoint-executable probe.

\begin{enumerate}

\item \textbf{Noise robustness of the coherence-collapse signature.}
Does the $C_{\ell_1}(k_{\rm peak})/C_{\ell_1}(0) \to 0$ trend survive
under realistic $T_1, T_2$ noise at levels achievable on current NISQ
hardware (superconducting ${\sim}100\,\mu$s $T_2$, ion-trap ${\sim}1\,$s
$T_2$)?
\emph{Basis:} This replication used pure statevector simulation. $C_{\ell_1}$
is by construction sensitive to off-diagonal decay, so $T_2$ dephasing
should attack it directly. The threshold at which the signature becomes
unresolvable is unknown from either the paper or this reproduction.
\emph{Next steps:} Add \texttt{qiskit\_aer.noise.NoiseModel} with
parameterized $T_1, T_2$ depolarizing/dephasing on each gate; sweep
$T_2 \in [1\,\mu\text{s}, 1\,\text{ms}]$ at $n=5,6,7$; recompute
$C_{\ell_1}(k_{\rm peak})/C_{\ell_1}(0)$ and identify the $T_2$
threshold below which the collapse ratio saturates. Compare against
published superconducting vs.\ ion-trap coherence numbers.

\item \textbf{Direct analog-vs-discrete correspondence.}
Does the Farhi--Gutmann analog Hamiltonian
$H = E(|w\rangle\langle w| + |s\rangle\langle s|)$ exhibit the same
coherence-collapse signature under continuous-time evolution,
iteration-for-iteration matched to the discrete algorithm reproduced
here?
\emph{Basis:} The paper is titled the ``discrete analogue of the analog''
algorithm; this replication only executed the discrete side. The
analog--discrete correspondence is asserted but not independently
verified here.
\emph{Next steps:} Implement Farhi--Gutmann $H$ in QuTiP or
\texttt{scipy.linalg.expm}; evolve $\rho(t)$ for $t \in [0, \pi/E]$;
sample $C_{\ell_1}(t), C_r(t), P(t)$ at Trotter step matched to
discrete $k$; overlay against discrete curves at $n=3\text{--}5$ to
test whether the two collapse profiles coincide.

\item \textbf{Extension to structured search.}
Does the coherence-collapse signature extend to structured search
problems (multi-marked-item Grover, fixed-point Grover, quantum-walk
search on graphs), or is it specific to canonical single-marked-item
Grover?
\emph{Basis:} The paper analyzes only canonical single-$|w\rangle$
Grover. Real applications (Grover-for-SAT, walk-based search) use
different oracle structures.
\emph{Next steps:} Extend \texttt{code/grover\_coherence.py} to
multi-marked ($M \in [1, N/4]$), fixed-point Grover (Yoder--Low--Chuang),
and quantum-walk-search on cycle/hypercube; compare
$C_{\ell_1}(k_{\rm peak})/C_{\ell_1}(0)$ trajectories to test
universality of the collapse.

\item \textbf{Hybrid analog--digital protocols.}
Could interleaved analog-Hamiltonian evolution and discrete
measurement/feedback retain more coherence while preserving the
$\sqrt{N}$ speedup? I.e., is Grover's coherence ``cost''
fundamental or protocol-dependent?
\emph{Basis:} The paper treats coherence as a monotone consumed during
search but does not distinguish Landauer-like fundamental cost from
protocol-specific cost.
\emph{Next steps:} Simulate a scheduled-weak-measurement Grover variant
(weak measurement at $k_{\rm opt}/2$, conditional continuation) and
compare $C_{\ell_1}$ trajectory vs.\ pure-unitary Grover at matched
final $P$. Also compare fixed-$E$ analog vs.\ time-dependent
$E(t)$ (Roland--Cerf adiabatic Grover).

\item \textbf{Mapping to real analog-platform hardware.}
How does the coherence-collapse signature manifest on Rydberg-atom
arrays (Lukin, Bernien), trapped-ion amplitude amplification (Monroe,
Blatt), and superconducting analog Grover on IBM/Google lattices, and
does the observed $T_2$-dependent degradation match theory?
\emph{Basis:} This is the experimental test of the paper's picture.
No public dataset was consulted or generated here.
\emph{Next steps:} (i) Literature search (Endres/Bernien 2017 Rydberg;
Figgatt 2017 ion Grover) for published $|\psi(t)\rangle$ tomography
during Grover-style protocols; (ii) if datasets exist, recompute
$C_{\ell_1}(t)$ on reconstructed $\rho(t)$ and compare with the
noise-augmented simulation from question~1; (iii) if none exist,
draft a minimal experimental proposal ($n=3$ Rydberg atoms, 4-shot
tomography every $100\,$ns).

\end{enumerate}
